%% =====================================================================================
\section{Limits of the method}\label{sec:limits}
%% =====================================================================================

The proof bounds the true loss $P_i(B_i)$ by a chain of steps: the union bound $\alpha_i\le\min(1,U_i)$ in the fibre
(\cref{lem:unionbound}); the pointwise hinge step, which discards the cap $\alpha_i\le1$ (\cref{lem:pointwise});
the comparison theorem, which aligns the progressions and tilts the measure (\cref{thm:comparison}); the enlargement
to all admissible moduli (\cref{lem:enlargement}); and the union bound over the levels in \cref{prop:F4}. On top of
this comes the slack of the numerical evaluation. In this section we discuss how much each step loses. Statements
labelled Proposition are proved. Statements labelled Observation are numerical findings from floating-point
experiments (Monte Carlo sampling at $m=16{,}000$, and exact computations and adversarial searches on small
analogues, and computations of the exact level bounds near $m=14{,}500$); they are not proofs and are not used
anywhere in the proofs of \cref{thm:main,thm:lean}.

\subsection{Steps that are sharp}

\emph{Alignment and tilt.} By \cref{prop:sharp}, \cref{thm:comparison} is attained by aligned cylinders and a
suitable $\nu$-bounded measure, so no better inequality holds for the whole class of $\nu$-bounded product measures.

\emph{Enlargement.} \Cref{lem:enlargement} loses nothing for a family that contains all admissible moduli
$m'p^j\ge m$ with $m'\mid Q_{i-1}$, up to the truncation of the exponents, which disappears as the exponents
$\gamma_j$ grow. Such families can be used at all levels at once, since moduli with different largest prime factors
are automatically distinct.

\emph{The redistribution rule.} By \cref{rem:rule-optimal}, given $P_{i-1}$ the rule \eqref{eq:rule} minimizes the
loss at level $i$ among all rules that preserve the marginal and multiply masses by at most $\nu_i$.

\emph{The tilt in practice.} \Cref{prop:sharp} uses a measure that is chosen freely, whereas $P_{i-1}$ is produced
by the rule \eqref{eq:rule}. The following observation shows that the tilt is nevertheless essentially attained.

\begin{observation}[the tilt is realized by nested shells]\label{obs:tilt}
Consider configurations of \emph{nested shells}, in which the progressions of every level are aligned at a common
centre, so that points that are heavily covered at one level are heavily covered at every smaller level as well. In
such a configuration the fibres of rich points receive, at each earlier level, the tilt
$\min(\nu_q,1/(1-\alpha_q))$ of \eqref{eq:rule}. Monte Carlo sampling ($10^6$ importance-sampled paths) at
$m=16{,}000$ with the final schedule, for $p\le2\cdot10^4$, gives the following totals of the per-level bound
$\E[(U_p-\delta_p)^+]/(1-\delta_p)$: $0.970$ with the full tilt $\nu_q$ (the bound used in the proof, whose exact value
is $0.9703$), $0.767$ without tilt ($\nu_q=1$), and $0.953$ with the tilt that nested shells realize. So the tilt
accounts for about $0.20$ of the budget, most of it at $200\le p<8{,}000$, and all but about $0.017$ of it is realized
by an actual configuration.
\end{observation}
% Source of the Observations: SCRATCH/agents/r9-tightness/FINDINGS.txt (sections 2, 3, 6 and the summary; the
% per-range Monte Carlo table there gives full tilt / no tilt / capped totals 0.97012 / 0.76699 / 0.75264 for
% p <= 2e4), SCRATCH/agents/r8-global-accounting/results.txt.  All numerical, floating point.

\subsection{The cap cannot be used by convex functionals}

The pointwise step replaces $(\min(1,U)-\delta)^+$ by the hinge $(U-t)^+$, which grows beyond $U=1$. The next
proposition shows that no convex functional can do better, so that the comparison theorem cannot exploit the cap
$\alpha\le1$.

\begin{proposition}\label{prop:cap}
Let $0<\delta<1$ and $f_\delta(u)=(\min(1,u)-\delta)^+/(1-\delta)$ for $u\ge0$. Let $\varphi\colon[0,\infty)\to\R$ be
convex with $\varphi\ge f_\delta$. Then for every random variable $W\ge0$ with $\E W<\infty$,
\[
  \E\varphi(W)\ge\min\Bigl(1,\ \min_{0\le t\le\delta}\frac{\E(W-t)^+}{1-t}\Bigr).
\]
\end{proposition}

\begin{proof}
Replacing $\varphi$ by its lower semicontinuous regularization, which is still convex, still majorizes the
continuous function $f_\delta$ and is not larger than $\varphi$, we may assume that $\varphi$ is continuous. Let
$\sigma^*=1/(1-\delta)$, the slope of $f_\delta$ on $[\delta,1]$, and let $\varphi'_-(1)$ be the left derivative at $1$.

\emph{Case $\varphi'_-(1)\le\sigma^*$.} Put $\sigma=\varphi'_-(1)$; it is a subgradient at the interior point $1$, so
$\varphi(u)\ge\varphi(1)+\sigma(u-1)\ge1+\sigma(u-1)$ for all $u\ge0$, because $\varphi(1)\ge f_\delta(1)=1$. Also
$\varphi\ge0$. If $\sigma\le0$, then $\varphi\ge1$ on $[0,1]$, and $\varphi\ge f_\delta=1$ on $[1,\infty)$, so
$\E\varphi(W)\ge1$. If $0<\sigma\le\sigma^*$, put $t=1-1/\sigma\le\delta$; then
$\varphi(u)\ge\max(0,\sigma(u-t))=\max(0,(u-t)/(1-t))$. For $t\ge0$ this gives $\E\varphi(W)\ge\E(W-t)^+/(1-t)$.
For $t<0$ it gives $\E\varphi(W)\ge(\E W-t)/(1-t)$, a convex combination of $\E W$ and $1$, which is at least
$\min(1,\E W)$, the bound for $t=0$.

\emph{Case $\varphi'_-(1)>\sigma^*$.} The continuous convex function $\psi(u)=\varphi(u)-\sigma^*u$ attains its
minimum over $[0,1]$ at some $u^*<1$, because $\psi'_-(1)>0$. It is then a global minimum on $[0,\infty)$: if
$\psi(v)<\psi(u^*)$ for some $v>1$, convexity would give $\psi(w)<\psi(u^*)$ for $w\in(u^*,1]$. Hence
$\varphi(u)\ge\varphi(u^*)+\sigma^*(u-u^*)$ for all $u\ge0$. Since
$\varphi(u^*)\ge f_\delta(u^*)\ge(u^*-\delta)/(1-\delta)$, this gives $\varphi(u)\ge(u-\delta)/(1-\delta)$, and together
with $\varphi\ge0$, $\varphi(u)\ge(u-\delta)^+/(1-\delta)$. So $\E\varphi(W)\ge\E(W-\delta)^+/(1-\delta)$.
\end{proof}

Applied to $W=U^{\mathrm{al}}$, \cref{prop:cap} says that the bound of \cref{thm:perprime}, minimized over $t$ and
capped at $1$, is the best that \cref{thm:comparison} can give for the loss. The cap does matter at the aligned
configuration: there the excess $\E[(U-1)^+]/(1-\delta)$ charged beyond the cap amounts to $0.217$ (numerical),
about $22\%$ of the budget for $p\le2\cdot10^4$, and between $20\%$ and $26\%$ of the budget of each range of primes.
An adversary, however, can spread the progressions apart.

\begin{observation}[spreading]\label{obs:spread}
In exact computations on small analogues (moduli built from the primes up to $7$ or up to $13$, with the worst
$\nu$-bounded measure and adversarial residues found by simulated annealing), spreading the progressions recovers
most of the gap between the capped aligned value and the per-level bound: between $59\%$ and $100\%$, depending on
the instance. At $m=16{,}000$, the explicit spreadings we tried recover at most about $15\%$ of the cap excess at the
prime $101$. Whether the true per-level worst case at this scale is much below the bound of \cref{thm:perprime} is
open.
\end{observation}

\subsection{Cross-level accounting}

The union bound over the levels in \cref{prop:F4} is not an identity, and its slack has a simple exact description.

\begin{proposition}\label{prop:live}
Call $x\in\Z_{Q_{i-1}}$ \emph{alive} if $x\notin\bigcup_{j<i}B_j$. Then
\[
  P_n\Bigl(\bigcup_{i=1}^nB_i\Bigr)=\sum_{i=1}^n\frac{\E_{P_{i-1}}\bigl[(\alpha_i-\delta_i)^+\,\one\{x\text{ alive}\}\bigr]}{1-\delta_i}.
\]
\end{proposition}

\begin{proof}
The union is the disjoint union of the sets $B_i\setminus\bigcup_{j<i}B_j$, which are $Q_i$-measurable, so by
\cref{lem:F1} its $P_n$-measure is $\sum_iP_i(B_i\setminus\bigcup_{j<i}B_j)$. The set $\bigcup_{j<i}B_j$ is
$Q_{i-1}$-measurable, so each fibre $F(x)$ either lies inside it or is disjoint from it. Hence
$P_i(B_i\setminus\bigcup_{j<i}B_j)=\sum_{x\text{ alive}}P_i(F(x)\cap B_i)$, and the computation in the proof of
\cref{lem:F3} finishes the proof.
\end{proof}

Thus $\sum_iP_i(B_i)$ also charges the losses on fibres that are already covered at an earlier level. In the nested
shells of \cref{obs:tilt} this happens systematically: the same rich points are charged at every level.

\begin{observation}[double counting]\label{obs:double}
On two small exact analogues (moduli dividing $2^43^25\cdot7$ and $2^53^35\cdot7$, with schedules optimized for the
bound of \cref{thm:perprime}), the per-level bounds sum to $1.0241$ and $0.9938$; simulated annealing finds systems
with $\sum_iP_i(B_i)=0.9934$ and $0.9511$, so the per-level chain loses only $0.031$ and $0.043$. Simulated annealing aimed
directly at $P_n(\bigcup_iB_i)$ finds at most $0.907$ and $0.886$: the double counting, about $0.07$ to $0.09$ of the budget, exceeds the
whole per-level slack. At $m=16{,}000$, in a Monte Carlo model of nested shells, the configuration that attains the
per-level bounds simultaneously (sum of capped losses about $0.74$ for $p\le2\cdot10^4$) newly covers only
$0.10$--$0.16$ of the mass. We did not obtain a rigorous improvement of the certified total from this observation.
\end{observation}

\subsection{The slack of the evaluation and the reach of the method}\label{sec:reach}

For the certificate of \cref{thm:lean} the evaluation itself is lossy. At $m=\MzeroLean$, with the schedule of that
certificate, the exact values of the level bounds (computed with \code{rigx.c}, without margins) total $0.982431$ over
the primes up to $\Pmax$, against the certified $\etaCertTwelve$. About $0.0097$ of the difference comes from
$50\le p<200$, where the decomposition of \cref{lem:AB} is inexact once $N_1>53$, and about $0.0053$ from $p>\PX$,
where the fractional-moment bounds are about six times the exact values. The exact evaluation of
\cref{sec:SBH,sec:tausplit,sec:lost} removes this slack. The following observation, from computations of the exact
level bounds without the safety margins and with heuristically optimized schedules, describes the reach of the
per-level method.

\begin{observation}[the frontier]\label{obs:frontier}
With the exact evaluation and the schedule of \cref{thm:lean}, the total crosses $1$ near $m=14{,}780$. With the
$\nKnotsMain$-knot schedule of \cref{thm:main} and $P_{\max}=10^9$, the total is $0.9999918$ at $m=14{,}500$ and
$1.0000041$ at $m=14{,}499$. One damped per-prime Newton step on the schedule (a schedule with $3{,}241$ knots)
gives $0.99996996$ at $m=14{,}498$ and $1.00002227$ at $m=14{,}495$. A total below $1$ is not
yet a proof: at $P_{\max}=10^9$ the test of \cref{prop:bbmst61} needs
$1-\eta\ge T_X/(\log k+\log\log k-3)^2k\approx4.5\cdot10^{-5}$, which is why $m=14{,}498$ to $14{,}500$ are not
certified. We estimate the frontier of the per-level method at $m^*\approx14{,}496\pm3$. Near $m=14{,}500$ the total is
a step function of $m$ that increases on average by about $0.015$ for each decrease of $1{,}000$ in $m$, and about
$\shareBelowTwoK\%$ of the budget is spent at the primes below $2{,}000$, where the level bound is now evaluated
exactly.
\end{observation}

A sharper evaluation of the same level bounds can therefore gain only a few units in $m$. Going substantially
further requires changing the method itself: a different rule for redistributing the removed mass, an accounting
across levels that exploits \cref{prop:live} (for instance a comparison theorem for the sub-stochastic kernels
$K_j\one\{\text{alive}\}$), or a non-convex treatment of the cap.
% Sources: the final report of SCRATCH/agents/r7-method-limit (frontier, 0.982431, 0.0097, 0.0053, 14,780, slope)
%          and certificate-14501/README.md (frontier data); share of p < 2000 from certificate-14501/output/dump.
% CHECK: sources disagree on the sensitivity of eta to m.  SCRATCH/agents/CONTEXT.md and r9-tightness/FINDINGS.txt
%        give "about 0.0085 per 1000" near m = 16,000; the r7 report measures about 0.015 per 1000 near m = 14,500
%        and says that 0.0085 does not hold there (r8-global-accounting/results.txt also gave about 0.015).
%        The text follows the r7 report.  (DELIV/README.md, "Limits of the approach", now also gives the limit
%        m ~ 14,500 and the frontier m* ~ 14,496, as here.)
